\section{Plano de ação resumido}
\label{sec:plano}

A Tabela~\ref{tab:plano} resume o plano nos quatro elementos exigidos pela orientação: estratégia, ações-chave, recursos necessários e indicadores de sucesso. As metas são hipóteses derivadas da Atividade IX e do desenho do piloto (Subseção~\ref{sub:piloto}) e deverão ser recalibradas com os dados observados. Os responsáveis seguem a estrutura organizacional da Atividade XI \cite{grupo2026a11}.

\begin{table}[htb]
\centering
\caption{Quadro-resumo do plano de ação}
\label{tab:plano}
\fontsizetable
\begin{tabularx}{\textwidth}{L{2.2cm} L{4.1cm} L{3.6cm} Y}
\toprule
\textbf{Estratégia} & \textbf{Ações-chave} & \textbf{Recursos necessários} & \textbf{Indicadores de sucesso (metas-hipótese)}\\
\midrule
\textbf{E1 -- Densidade local por parcerias (principal)} &
1. Descoberta e piloto assistido de oito semanas: 15 a 20 entrevistas com tutores, cinco clínicas e duas organizações; 10 testes de uso sem ajuda; 40 a 50 tutores e 2 a 3 clínicas no piloto (Giovanna, Júlia).\newline
2. Fechar três clínicas (três meses gratuitos do Plano Parceiro) e duas ONGs ou lideranças comunitárias (Giovanna).\newline
3. Instalar QR e \textit{links} por parceiro e evento; divulgação em redes sociais e indicação com créditos (Júlia).\newline
4. Corrigir segurança e privacidade antes do piloto com dados reais (Gabriel, Kathleen, Júlia).\newline
5. Medir CAC, ativação e retenção; aplicar os gatilhos da Tabela~\ref{tab:gatilhos} antes de abrir nova praça. &
Tempo dos quatro integrantes, com registro das horas mesmo sem remuneração; materiais de divulgação (R\$~150,00 mensais já previstos) e teto de aquisição (cerca de R\$~438 mensais); consultoria veterinária (R\$~600,00 mensais) e jurídico e LGPD (R\$~250,00 mensais), já nos custos fixos; infraestrutura em nuvem (R\$~100,00 mensais). &
Usuários ativos: 300 (mês 3, condicionado à decisão da semana 8), 800 (mês 6), 1.500 (mês 12).\newline
Clínicas ativas: 3, 6 e 10 nesses marcos.\newline
Ativação: 60\% dos inscritos no piloto.\newline
Retenção no dia 30: 30\% dos inscritos.\newline
CAC observado $\le$ teto (R\$~3,29).\newline
Atendimentos intermediados: 10\% dos usuários ativos por mês.\\
\midrule
\textbf{E2 -- Aprofundamento de receita e novos serviços (complementar)} &
1. Aplicar o teste de aceitação do preço da Atividade IX, perguntando às clínicas o que valorizariam e qual faixa de preço considerariam (Júlia, Giovanna).\newline
2. Lançar o PetCuida+ com plano anual e meta de conversão de 5\% a 8\% (Giovanna, Gabriel).\newline
3. Vender campanhas patrocinadas e ao menos um convênio formal para financiar o fundo (Giovanna).\newline
4. Testar a intermediação somente após definir repasse, transparência e requisitos (colegiado).\newline
5. Priorizar novos serviços da lista de espera somente com demanda comprovada (colegiado). &
Tempo de desenvolvimento (Gabriel e Kathleen); taxa de conta de desenvolvedor da loja de aplicativos; material comercial para patrocinadores; parecer jurídico para contratos e termos. &
Conversão ao PetCuida+: $\ge$ 5\% (meta 8\%).\newline
Fundo de liquidação por atendimento: de R\$~6,60 para R\$~3,30.\newline
Receita recorrente: $\ge$ 40\% da receita (sem convênios).\newline
Convênio ou patrocínio formal: $\ge$ 1.\newline
Duas clínicas dispostas a repetir um ciclo controlado.\newline
Contribuição por usuário: de R\$~0,55 para R\$~1,42.\\
\bottomrule
\end{tabularx}
\fonteautores
\end{table}

\subsection{Definições dos indicadores de aprendizagem}

Curtidas e \textit{downloads}, isoladamente, não demonstram valor recorrente. Por isso, o piloto adota definições operacionais explícitas, apresentadas na Tabela~\ref{tab:indicadores}. Caso a amostra seja menor que a prevista, o grupo reporta as contagens absolutas e as limitações, sem extrapolar percentuais.

\begin{table}[htb]
\centering
\caption{Indicadores de aprendizagem do piloto}
\label{tab:indicadores}
\fontsizetable
\begin{tabularx}{\textwidth}{L{2.4cm} Y L{4.2cm}}
\toprule
\textbf{Frente} & \textbf{Definição} & \textbf{Meta inicial proposta}\\
\midrule
Descoberta & Entrevistas e interlocutores locais. & 15 a 20 tutores; 5 clínicas; 2 organizações.\\
Usabilidade & O participante conclui as tarefas-chave sem ajuda. & 10 testes; ao menos 8 concluem; 80\% explicam QR e créditos.\\
Ativação & Perfil do animal completo e primeiro lembrete cadastrado em até sete dias. & 60\% dos inscritos no piloto.\\
Retenção & Ao menos uma ação relevante entre os dias 22 e 30 (consultar o histórico, atualizar um cuidado preventivo ou iniciar um encaminhamento). & 30\% dos inscritos no piloto.\\
Oferta & Parceiro disposto a repetir um ciclo controlado. & 2 clínicas interessadas em novo teste.\\
\bottomrule
\end{tabularx}
\fonteautores
\end{table}

\subsection{Cronograma e marcos}

\begin{alineas}
  \item \textbf{Meses 0 a 3 (preparação e piloto):} descoberta, revisão e testes (semanas 1 a 4); correções de segurança e privacidade; constituição da empresa; teste de aceitação do preço; piloto assistido com 40 a 50 tutores e duas a três clínicas (semanas 5 a 8); decisão da semana 8; três clínicas, duas ONGs e meta de 300 usuários ativos ao final do mês 3, se a decisão for favorável;
  \item \textbf{meses 4 a 6 (validação do núcleo):} medição de CAC, retenção e atendimentos por usuário; seis clínicas e 800 usuários; revisão do fundo de liquidação;
  \item \textbf{meses 7 a 12 (consolidação):} dez clínicas, 1.500 usuários, 75 assinantes e ao menos um convênio; decisão colegiada sobre abertura de nova praça;
  \item \textbf{meses 13 a 24 (expansão condicionada):} abertura de uma nova área da região metropolitana se os gatilhos forem atingidos; meta de 5.000 usuários no mês 24; avaliação de financiamento externo, pois o cenário com equipe remunerada depende de escala maior ou de recursos de terceiros.
\end{alineas}

\subsection{Governança e papéis no próximo ciclo}

A execução segue as regras da Atividade XI: Giovanna aprova escopo, custos e parcerias e coordena as entrevistas e os parceiros; Gabriel responde pela arquitetura e pela segurança técnica; Júlia, pela privacidade, pela interface, pela comunicação externa e pelos testes de usabilidade; e Kathleen, pela qualidade, pela liberação de versões e pela ponte de comunicação entre as frentes, sintetizando dúvidas e decisões. Essa divisão é uma proposta, sujeita à disponibilidade real dos integrantes. A revisão mensal de negócio e de carga compara os indicadores das Tabelas~\ref{tab:plano} e~\ref{tab:indicadores} com as metas e decide se o plano deve ser mantido, ajustado ou interrompido.
